\section{All powers of the corrected phase detectors}
\label{sec:actual-phase-detector-moments}

The rank comparisons must survive passage to a phase. A bound on the
unselected one- and two-point functions is insufficient: we need bounds
on every power of the complete corrected operators. We now obtain those
bounds from the same finite-cutoff corrected root estimates.

Retain the hierarchy, its reflection-positive Hilbert space, the bounded
slab operators, and their ground-space compressions $M(w)$ from
Section~\ref{ex-phases}. We use their permutation, empty-word and adjoint
laws, the separated-cluster limit of Lemma~\ref{ex-cluster-train}, and
the reflected-word identification established by the real phase-grid
comparison. In particular, if $\bar w$ denotes ordinary conjugation of
all labels without reflecting its points, then
\begin{equation}\label{eq:episode-conjugate}
                         M(\bar w)=M(w)^*.
\end{equation}
This last identity is a substantive input: reverse the order using
permutation invariance, apply the reflected-word identification, and
then the slab-adjoint law. The real phase-grid producer must supply that
identification on this same hierarchy.

Fix a grade and one of the comparison trials of
Section~\ref{sec:actual-rank-purchases}. Write $k$ for its grid depth,
$b\leq k$ for its root-test depth, and $g\geq0$ for its corrected root
price. At its finitely many row sites $x_j$ let $f_j^{N,p}$ be the
actual cutoff insertion with the trial's lifted label. Along the
selected mesh-then-period order those labels converge to labels $f_j$
on the same hierarchy. Write $a_j$ for the prescribed free one-point
value, and $\kappa_{jh}$ for the prescribed connected pair value.
The trial supplies chamber interiority and, for every $\varepsilon>0$,
eventually in that order,
\begin{align}
 \|f_j^{N,p}-a_j\|_{\mathrm{root},b}&\leq g+\varepsilon,
                                                        \label{eq:episode-root-one}\\
 \|(f_j^{N,p}-a_j)(f_h^{N,p}-a_h)-\kappa_{jh}\|_{\mathrm{root},b}
       &\leq g+\varepsilon\quad(j<h).                   \label{eq:episode-root-two}
\end{align}
These are the corrected polynomial estimates purchased before extraction,
not estimates of a product of uncorrected root norms. Constant functions,
sums and products preserve the chamber-interior property, so it also
holds for the displayed corrections.

\subsection{A finite polynomial with coupled cutoff and limiting labels}

Consider a finite sum of words
\[
 Z^{N,p}(u)=\sum_{\alpha=1}^L c_\alpha
     \prod_{j=1}^{n_\alpha}
       f_{\alpha j}^{N,p}(x_{\alpha j}+u),
 \qquad Z=\sum_{\alpha=1}^L c_\alpha M(w_\alpha).
\]
Each word uses sites on the depth-$k$ grid that are generic: in every
coordinate its distinct slots have distinct coordinates. Different words
in the sum may share sites. Empty words are allowed and represent the
constant $1$ at both cutoff and limit. The coefficients $c_\alpha$ are
fixed before the cutoff limits. Each cutoff label converges in the
selected order to the label of the corresponding slot of $w_\alpha$.
We call this coupled finite list a polynomial packet. Suppose its
cutoff polynomial at $u=0$ is chamber-interior and has root price at most
$g+\varepsilon$ at depth $b$, eventually for every $\varepsilon>0$.

For $q\in\mathbb N$ and integers $t\geq1$ place $2q$ translates of this whole polynomial at
physical displacements
\[
                   u_j(t)=2ta_bj(1,1,1,1),\qquad 0\leq j<2q.
\]
These are depth-$k$ grid displacements because $b\leq k$ and the scales
are aligned. All coarse cell addresses are even in every coordinate.
Consequently copying from the home cell to them is ordinary translation,
without a reflection or conjugation of the copied polynomial. At even
$j$ use $\overline{Z^{N,p}}$, and at odd $j$ use $Z^{N,p}$.
Conjugation does not change a root price: it conjugates the full copied
product and its expectation, preserving its absolute value.

For each fixed $t,q$, sufficiently large periods make the $2q$ cell
addresses distinct. Indeed $2t(2q)$ not exceeding the coarse period
suffices; reducing modulo that period is then injective on the chosen
addresses. The planted inequality on the same Wilson law gives
\begin{equation}\label{eq:episode-cutoff-power}
 \left|\mathbb E_{N,p}
       \prod_{j=0}^{2q-1}
       \begin{cases}
        \overline{Z^{N,p}(u_j(t))},&j\ \text{even},\\
        Z^{N,p}(u_j(t)),&j\ \text{odd}
       \end{cases}\right|
                    \leq(g+\varepsilon)^{2q}.
\end{equation}
The assertion is eventual in period and then in the mesh tail at that
period. It makes no claim at periods too small to hold the distinct
copies. For $q=0$ both sides equal $1$.

Expand the product into its finitely many choices of a monomial for
each of the $2q$ packets. Within each chosen monomial the slots are
generic. Between slots in different packets, equality of any one
coordinate is an affine equation in $t$ with nonzero slope, since their
packet indices differ. It excludes at most one value of $t$. There
are finitely many slot pairs, coordinates and monomial choices. Thus
all sufficiently large $t$ make every expanded configuration generic.
For each such fixed $t$, take the fine-mesh and period limits in their
prescribed order. The same limiting labels are used in every term.
Ordinary conjugation of a cutoff insertion agrees with conjugation of
its label once its fixed reference depth is admitted. Convergence of
these conjugated labels follows from continuity of conjugation.
The point-limit theorem therefore applies term by term. At finite
cutoff the integrands are continuous on the compact configuration
space, so their finite sums are integrable and expansion under the
integral is legitimate. This proves a limiting finite sum $P_q(t)$
with
\begin{equation}\label{eq:episode-first-limit}
                    |P_q(t)|\leq(g+\varepsilon)^{2q}
\end{equation}
for all sufficiently large $t$.

Now let $t\to\infty$. The time-coordinate separation between packets
$j<h$ is $2ta_b(h-j)\to\infty$. The separated-cluster theorem permits
the simultaneous spatial translations used here. Apply it to each
expanded monomial, retaining the order of its packets. By
\eqref{eq:episode-conjugate}, the even packets converge to adjoints of
their compressed whole-word operators. Summing the finitely many
terms gives
\[
       P_q(t)\longrightarrow
       \langle\Omega_0,(Z^*Z)^q\Omega_0\rangle.
\]
This is simply expansion of the ordered product
$Z^*Z\cdots Z^*Z$; commutativity is not needed for this identification.
Take real parts in \eqref{eq:episode-first-limit} and then let
$\varepsilon\downarrow0$. We have proved
\begin{equation}\label{eq:episode-all-powers}
  \langle\Omega_0,(Z^*Z)^q\Omega_0\rangle\leq g^{2q}
                    \qquad(q=0,1,2,\ldots).
\end{equation}
The left side is real and nonnegative. The order of limits was:
fix $q,\varepsilon$; fix sufficiently large $t$; take the cutoff
limits; take $t\to\infty$; finally take $\varepsilon\downarrow0$.
No uniform-in-$q$ cutoff threshold or interchange of those limits
has been used.

\subsection{The two actual corrected polynomials}

Apply this result first to the two-term packet with coefficients $1,-a_j$
and words $[f_j],\varnothing$. Its cutoff polynomial is exactly
$f_j^{N,p}-a_j$, and its compressed limit is
\[
                           D_j=M(f_j)-a_jI.
\]
Use the actual trial's lifted labels at cutoff and its limiting labels
in the compressed word. Genericity is inherited from the trial's row.
Thus \eqref{eq:episode-root-one} supplies the required root bound.

For $j<h$ use the four words
$[f_j,f_h],[f_h],[f_j],\varnothing$ with coefficients
$1,-a_j,-a_h,a_ha_j-\kappa_{jh}$. The finite-cutoff polynomial is
\[
 f_j^{N,p}f_h^{N,p}-a_j f_h^{N,p}-a_h f_j^{N,p}
                  +a_ha_j-\kappa_{jh}
       =(f_j^{N,p}-a_j)(f_h^{N,p}-a_h)-\kappa_{jh}.
\]
Its compressed limit is the entire operator
\[
 D_{jh}=M(f_jf_h)-a_jM(f_h)-a_hM(f_j)
                         +(a_ha_j-\kappa_{jh})I.
\]
The two-slot monomial uses the actual distinct row slots $j,h$ even
when their labels coincide or one is the unit. Its genericity follows
from that of the row; no removal or renaming of slots is required.
Equation~\eqref{eq:episode-root-two} therefore applies. We obtain
\eqref{eq:episode-all-powers} for both $D_j$ and $D_{jh}$ with the
same $g$.

The ground algebra has the full-support spectral measure constructed
in Section~\ref{ex-phases}. Lemma~\ref{ex-character-moments} now gives,
for every character $\chi$,
\[
              |\chi(D_j)|\leq g,\qquad |\chi(D_{jh})|\leq g.
\]
Hence its phase correlations obey
\[
 |S_\chi(f_j)-a_j|\leq g,\qquad
 |S_\chi(f_jf_h)-(\kappa_{jh}+a_ja_h)|
                  \leq g(1+|a_j|+|a_h|).
\]
These estimates hold at every comparison trial and for every character,
including characters in null sets for the representing measure. Along
the purchased trials $g\to0$ and the one-point comparison values are
bounded, by continuity of the free correlations at the limiting test
configurations. The finite separating tests and the eventual spanning
property of the same limiting lifts are those retained in the rank
construction. The finite-dimensional argument of
Proposition~\ref{ex-no-new-null} consequently applies in each grade:
for every represented phase whose correlations respect the original
null relations, its null space equals that of the original hierarchy.
This proves the detector-moment input to phase retention, conditional
on the stated slab, reflected-word and same-trial inputs; it does not
construct a missing reflected-word identification.
